\section{Compact coordinates for certified component inventories}
\label{sec:compact-coordinates}

The chapter in
\code{ProveIt\_UnknotRecognition\_Research\_2026-10-09.zip} gives a useful
bound on the number of distinct connected normal vectors in one supplied
surface. Its matching-space argument also identifies redundant coordinates
in our weighted census. We now use this observation to reduce the weights
transported by the native \code{coordinates} query, while returning exactly
the same full-coordinate inventory. This section separates that local
implementation improvement from the unresolved search for useful surfaces.

\subsection{A linear bound on the number of component types}

Let $F$ be an embedded normal surface in a finite compact triangulated
3-manifold with $t$ tetrahedra and $v$ quotient vertices, whose links are
connected spheres or discs. Let $Q$ be the set of positive quadrilateral
coordinates of $F$, and put $q=|Q|\le t$. A \emph{type} means a distinct
full $7t$-coordinate vector of a connected component, with repeated copies
stored by a binary multiplicity. The report proves
\[
 H\le v+2q,\qquad
 H\le v+q\quad\hbox{if all components are two-sided}.
\]
These statements concern one simultaneously embedded surface. They do not
bound the number of candidate normal surfaces in the triangulation.

Here is a proof, including the role of one-sided components. In the real
matching space $V_Q$ supported on $Q$, the kernel of projection to the $q$
quadrilateral coordinates consists exactly of the $v$ vertex-link directions.
Indeed, triangle-only face equations identify the corner coefficients around
each quotient vertex, whose corner adjacency graph is connected. Thus
$\dim V_Q\le v+q$.

Distinct full vectors of disjoint connected two-sided normal surfaces are
linearly independent. Otherwise, clearing denominators in an integer-matrix
kernel gives an equality between positive integer combinations on disjoint
sets of types. Take the required parallel copies in disjoint normal product
neighborhoods. Equality of full normal coordinates implies normal isotopy,
which preserves the multiset of connected vectors. This contradicts the
disjoint sets of types. The coordinate uniqueness used here is the classical
normal-surface statement recorded in Burton's Lemma~2.5\cite{burton-convert}.

Choose one representative of each nonvertex type of $F$. Replace each
one-sided representative by the horizontal frontier of a small normal
interval neighborhood. Its nontrivial normal line bundle makes this a
connected two-sided surface with twice the original vector. For properly
embedded surfaces with boundary, the vertical boundary lies in the ambient
boundary. Retain two-sided representatives and coalesce equal resulting
vectors. These $r$ representatives can be taken mutually disjoint and
disjoint from all $v$ small vertex links. Independence gives
$r+v\le\dim V_Q\le v+q$, hence $r\le q$. Each resulting vector $z$ has at
most two original preimages, $z$ and $z/2$. Adding at most $v$ triangle-only
types proves $H\le v+2q$. When all components are two-sided the replacement
is injective and gives the stronger bound.

The factor two is necessary. In the one-tetrahedron layered solid torus,
identify face $0$ with face $1$ by the cycle $(0\,1\,2\,3)$ and leave the
other faces on the boundary. In our coordinate order, the vertex-link disc
and a M\"obius band have vectors
\[
 \ell=(1,1,1,1;0,0,0),\qquad m=(0,0,0,0;0,1,0).
\]
The source $\ell+3m$ decomposes into the three types $\ell,m,2m$, attaining
$v+2q=3$. The source $\ell+2m$ attains the two-sided bound $v+q=2$.
Our native tests preserve the distinction between the M\"obius band and its
annular double, including large binary multiplicities.

If source coordinates have at most $B\ge1$ bits, each component coordinate
is bounded by its source entry, and multiplicities are bounded by the total
number of normal discs. Since $v\le4t$, the coalesced full inventory requires
$O(t^2(B+\log(t+1)))$ bits. AHT already provide polynomial-time component
recovery from binary normal coordinates\cite[Corollary 17]{aht-orbits}.
The present argument sharpens output accounting; it is not a new proof of
polynomial-time normal-surface classification. Equality of the sum of
proposed component vectors with the source is insufficient evidence:
normal addition can reconnect surfaces. Orbit membership must still be
certified.

\subsection{An injective additive coordinate projection}

For each quotient vertex choose a triangle coordinate $a_V$ minimizing its
value in the supplied source $x$, with ties broken by flat coordinate index.
Retain every positive source quadrilateral coordinate and each anchor with
$x_{a_V}>0$. Let $a$ be the number of retained anchors. The transported
dimension is now
\[
 d=\max(1,q+a),\qquad a\le v,
\]
where the single dummy weight handles the empty source under the existing
weighted-orbit API. All omitted source-zero quadrilaterals and anchors are
zero on every component, because component vectors are nonnegative and sum
to $x$. Selecting the source minimum maximizes the opportunity to omit an
anchor: a vertex with any zero triangle needs no anchor weight.

The projection is injective on possible component vectors. Two vectors with
the same retained entries have equal quadrilaterals; their difference is
triangle-only, so its corner entries are constant around each vertex. Its
anchor value is zero at every vertex, forcing the difference to vanish.
Consequently equal projected orbit weights mean equal full component vectors.
Coalescing by projected weight loses no component types or multiplicities.
This is a source-dependent linear projection, not a claim that quadrilateral
coordinates alone distinguish vertex links.

For a glued face, matching has the form
\[
 T_{t,v}+Q_{t,\alpha}=T_{u,p(v)}+Q_{u,\beta}.
\]
After assigning quadrilaterals and anchors, propagate
$T_{u,p(v)}=T_{t,v}+Q_{t,\alpha}-Q_{u,\beta}$ along a corner spanning forest.
The producer constructs this forest once from the face pairings. Each
histogram entry then takes $O(t)$ integer additions and comparisons to
recover, followed by the existing full-vector matching and topology checks.
Values are checked against the original source. With $H=O(t)$, decoding
adds $O(t^2)$ arithmetic operations on $O(B+\log(t+1))$-bit integers;
source indexing and sorting have their usual additional polynomial costs.

The orbit relation and its trace are unchanged. The existing disk-corner
weight construction now selects only the retained coordinates, reducing
both dense vector dimension and the number of initial indicator intervals.
The number of selected indicators is at most $q+a$. Weighted transport and
its independent prefix-sum arithmetic are unchanged. For a meridian in our
$t$-tetrahedron layered fixture, $q=t$ and $a=0$, so $d$ drops from $7t$ to
$t$. In general this improves a local polynomial routine; it does not change
the asymptotic complexity class of the complete recognizer.

\subsection{Certificate compatibility and independent reconstruction}

New full-coordinate censuses use \code{normal-component-census-v2}.
The source fingerprint, boundary basis and checked weighted-orbit trace
remain mandatory. The selected indices are recomputed from the original
source; no claimed coordinate basis or spanning tree is trusted. The
summary still contains full $7t$-coordinate vectors. The public
\code{weight\_dimension} now reports the smaller transported dimension.
Disk and summary queries continue to emit version one with three and five
weights respectively. The verifier retains the original dense interpretation
of all version-one coordinate certificates.

The checker selects anchors independently. Instead of the producer's
face-pair spanning forest, it builds an adjacency index of the validated
matching equations and repeatedly solves a row with one unknown triangle.
It checks source bounds and then runs the full matching/signature validation
on every recovered vector. Its arithmetic decoder neither imports nor calls
the producer's coordinate-basis module. Shared finite-manifold geometry
validation remains part of the pre-existing trust contract. The weighted
checker still performs its separate transport replay.

Tests exercise both schemas, reject changing a dense proof's version or a
compact proof's version, disable producer reconstruction during verification,
and alter both projected and recovered coordinates. Other cases include
empty surfaces, several quotient vertices, relabelings, zero anchors,
triangle-only sources, cancellation late in reconstruction, and a multiplicity
$2^{5000}+1$ without sheet expansion. The finite compact torus-boundary
contract remains unchanged, as does the separate obligation to certify a
knot-diagram-to-triangulation correspondence.

\subsection{What the bound leaves open}

The type bound controls the size of one completed inventory. It does not
control how many quadrilateral supports a recognizer must try, how encoded
coordinates grow under a geometric hierarchy, or the size of triangulations
after cutting and simplification. Nor does a histogram of full component
vectors label individual boundary ports for arbitrary future gluing. Those
operations still need the incidence information and sufficiency arguments
discussed in the sparse-observer sections. Reducing an existing polynomial
component query is useful implementation progress, but supplies none of those
missing global search or state bounds by itself.

\subsection{Native validation and paired measurements}

The pinned baseline is \code{aeb891e7159a}. The driver
\code{weighted\_research/compact\_coordinates.py} records 442 current source
and fixture hashes, separately records every baseline package hash, and checks
that the source hashes do not change during a run. The dense baseline is
loaded in an isolated package namespace from Git; it includes the same
adaptive interval-merger implementation. The only changed pre-existing
runtime modules are the component geometry, producer and verifier; the
compact producer decoder is a new module.

The frozen Regina 7.4 corpus contains 1,275 vectors over 45 triangulations.
All three census modes and the normalized disc-count query pass 5,100 oracle
comparisons and independent certificate replays. The extra compact audit
checks full-summary equality against the dense predecessor on every vector,
replays all 1,275 dense certificates with the current verifier, and confirms
that every unweighted orbit proof is unchanged. The source-support type
bounds pass on every corpus entry, including 252 inputs with nonorientable
components. These are finite audit checks accompanying the proof above.
The audit takes 35.548 seconds, including the additional predecessor runs;
it overlaps the test suite and is not a performance measurement.

The maintained suite passes all 1,125 tests in 220.141 seconds. After tests,
audit, commit and publication finished, the benchmark ran serially for
267.924 seconds. It uses six supplied surfaces, four shuffled arms (dense,
compact, and an A/A copy of each), five measured rounds and one retained
warm-up round. All 120 measured calls and 24 warm-ups complete successfully.
Every timed call includes fresh source validation, coordinate-weight
construction, orbit discovery, certificate production and independent replay.
Fixture creation, serialization and cross-arm equality checks are outside
the timers. Both arms return identical full component summaries. Raw samples,
source pins and deduplicated full certificates are retained in
\code{data/compact-coordinates-benchmark.json}.

\begin{center}\small
\begin{tabular}{lrrrrr}
\toprule
Input & Dense ms & Compact ms & Dense/compact & Dense A/A & Compact A/A \\
\midrule
meridian-1 & 1.003 & 0.961 & 1.035 & 0.986 & 1.009 \\
meridian-16 & 54.962 & 22.945 & 2.413 & 0.993 & 0.994 \\
meridian-64 & 2204.769 & 484.808 & 4.548 & 0.988 & 1.042 \\
meridian-128 & 15969.781 & 2744.039 & 5.851 & 0.991 & 0.993 \\
triangle-only & 5.332 & 4.005 & 1.331 & 0.988 & 0.993 \\
binary-1024 & 28.181 & 16.203 & 1.705 & 1.014 & 0.996 \\
\bottomrule
\end{tabular}
\end{center}


At 128 tetrahedra, dimension falls from 896 to 128, and the initial and peak
piecewise-constant weight-run count falls from 385 to 257. The underlying
133 orbit cycles and 21,008 replay events are unchanged. Median total time
falls from 15.970 seconds to 2.744 seconds; the median within-round ratio is
$5.851$, which need not equal the ratio of separate medians. Certificate
size falls from 2,504,362 to 1,867,476 bytes in the recorded compact JSON
encoding. This is a complete full-coordinate-query comparison, not a
comparison of different output modes.

The corresponding paired gains are $2.413$ at 16 tetrahedra and $4.548$ at
64. Dense A/A ratios range from $0.986$ to $1.014$ and compact A/A ratios
from $0.993$ to $1.042$. The one-tetrahedron $1.035$ ratio is small enough
that it should not be promoted as a meaningful gain. A triangle-only source
with two vertex-link types improves by $1.331$; an eight-tetrahedron meridian
scaled by $2^{1024}$ and supplemented with $2^{1024}+1$ vertex links improves
by $1.705$. These are six supplied-source cases, not a universal speedup
claim or new automatic unknot-recognition coverage.

